Este capítulo retoma o percurso do trabalho em três movimentos. O primeiro verifica, objetivo por objetivo, o que foi efetivamente entregue. O segundo reúne os resultados numéricos que sustentam as escolhas de projeto. O terceiro delimita o que ficou fora do alcance e indica o que a continuidade do trabalho exigiria.

\section{Alcance dos objetivos}\label{sec:5.1}

O primeiro objetivo específico consistia em levantar, na literatura, os parâmetros arquiteturais de uma memória cache e justificar as escolhas adotadas. A seção~\ref{sec:2.5.1} respondeu à parte descritiva, e a seção~\ref{sec:3.2} fixou a configuração do projeto em 8 KiB por cache, linha de 16 bytes, quatro vias e política LRU. No entanto, a justificativa apoiava-se em duas fontes de terceiros: a arquitetura de referência CVA6 fornecia os valores, e os dados de síntese do BRISC-V indicavam o custo em área de cada parâmetro. Assim, faltava o comportamento do próprio subsistema sob carga, lacuna essa preenchida no \refCap{4_Verificacao}, onde a varredura mostra que 8 KiB e quatro vias se situam exatamente no ponto em que o ganho marginal se anula. Por conseguinte, o objetivo está alcançado com uma justificativa mais forte do que a inicialmente prevista.

O segundo objetivo pedia a especificação da máquina de estados que governa acertos e faltas. O \refCap{3_Implementacao} a fixou em oito estados, com as condições de transição e os protocolos de troca com o pipeline e com a memória descritos em termos de sinais e de temporização. A separação entre o caminho de acesso e o caminho de consistência, com cinco e três estados respectivamente, decorreu dessa especificação. O objetivo está alcançado.

O terceiro objetivo tratava da implementação das duas caches como módulos distintos, com a observação de que a cache de instrução dispensa o bit de sujeira. A seção~\ref{sec:3.1} descreve a solução adotada: uma única classe, instanciada duas vezes com um parâmetro de elaboração que suprime o caminho de escrita na primeira. O objetivo está alcançado.

O quarto objetivo pedia o controlador de faltas, responsável por suspender o pipeline, sinalizar a espera ao estágio consumidor e retomar a execução. As seções~\ref{sec:3.5} e~\ref{sec:3.7} descrevem a integração e a máquina de estados. O objetivo está alcançado.

O quinto objetivo tratava do mecanismo de consistência exigido pela instrução FENCE.I. A seção~\ref{sec:3.9} descreve as duas máquinas que o implementam, uma no controlador de cada cache (os estados de drenagem) e outra no coordenador, e a seção~\ref{sec:4.3} apresenta a verificação. O teste de código automodificável é o que dá sustentação ao resultado, já que a sub-rotina é reescrita, a sincronização é executada e a segunda chamada devolve o valor novo. Como esse desfecho exige tanto a drenagem das linhas sujas quanto a invalidação da cache de instruções, qualquer das duas omissões faria o teste falhar. O objetivo está alcançado.

O sexto objetivo pedia a verificação por testes automatizados sobre um modelo de referência. O \refCap{4_Verificacao} reporta trinta suítes e cento e oitenta casos, dos quais cento e setenta e nove aprovados, nenhum reprovado e um cancelado por ausência de variável de ambiente. Mais relevante que o número é o desenho, que os casos de varredura comparam não apenas os dados, mas os contadores, contra um modelo escrito independentemente em Scala. O objetivo está alcançado.

O sétimo objetivo pedia quantificar o efeito da geometria sobre a taxa de acerto e sobre o tráfego de memória em bytes, além de determinar o custo em ciclos de cada evento do subsistema. As três partes estão cumpridas. As duas primeiras, pela varredura de treze configurações sobre um trace fixo. A terceira, pela instrumentação do custo da instrução de sincronização, que resultou em vinte medições exatas na cache isolada e dezesseis no pipeline. O objetivo está integralmente alcançado.

\section{Síntese dos resultados}\label{sec:5.2}

A configuração adotada obteve taxa de acerto de 69,5\% sobre o trace de avaliação, com 457 faltas e 16 despejos em mil e quinhentos acessos. Isoladamente, esse número não significa nada, o que o qualifica é a varredura que o cerca.

Nos eixos de capacidade e de associatividade, os incrementos de desempenho se anulam exatamente na configuração escolhida. Elevar a capacidade de 1 KiB para 8 KiB produz ganhos sucessivos de 6,3, 7,2 e 4,1 pontos percentuais; dobrar para 16 KiB acrescenta 0,1 ponto, ou um único acerto em mil e quinhentos acessos. O mesmo padrão se repete na associatividade, com incrementos de 1,7, 0,6 e 0,1 ponto ao passar de uma para duas, quatro e oito vias. O resultado tem duas leituras. A primeira é que a geometria adotada não é arbitrária: ela está no ponto de saturação. A segunda é que o efeito mais informativo não está na taxa de acerto, e sim nos despejos, que caem de 261 para 16 ao longo do eixo de capacidade e de 65 para 8 ao longo do de associatividade. A cache de base já captura quase todas as linhas que seriam despejadas sujas, e é esse fato, mais que o acerto, que sustenta a escolha diante do custo em área.

O eixo de tamanho de linha é o único em que a configuração adotada não é a de melhor taxa de acerto, e por isso o mais informativo. A linha de 64 bytes acerta 84,2\% contra 69,5\% da de 16 bytes, uma diferença de 14,7 pontos percentuais e 220 faltas. Lido apenas por esse critério, o resultado condenaria a escolha. Lido em bytes transferidos, inverte: a linha de 64 move 19.200 bytes contra 7.568 da de 16, ou 2,5 vezes mais, para evitar 220 faltas. A conclusão que se extrai não é que uma das duas esteja errada, mas que a taxa de acerto é um critério insuficiente quando o custo da transferência é proporcional ao seu volume, e que a escolha de 16 bytes se justifica por tráfego e por área, não por acerto.

\section{Trabalhos futuros}\label{sec:5.3}

A continuidade mais direta deste trabalho está em duas etapas de custo baixo que fecham lacunas declaradas. A primeira é um modo de contorno das caches no pipeline, que tornaria possível a comparação entre execução com e sem cache e converteria o custo por evento, já medido, em ganho agregado de desempenho. A segunda é a execução do projeto na placa, com leitura dos contadores de desempenho em hardware e confirmação da pré-carga do programa na memória principal, que fecharia a distância entre o que a síntese reporta e o que o subsistema faz em operação.

Uma segunda frente decorre do que a análise de tráfego revelou. Se o custo da transferência for proporcional ao volume, e não ao número de transações, a linha de 16 bytes é a escolha correta e a de 64 não se justifica. Se o modelo de memória for alterado para refletir largura de banda, o compromisso se desloca, e vale reexaminá-lo.

Uma terceira frente diz respeito ao próprio controlador. A cache implementada é bloqueante, e o estado da arte em cache de dados de código aberto~\cite{million2024hpdcache} opera com múltiplas faltas em voo. A implementação de uma cache não bloqueante, com registro de faltas em andamento, é o passo natural de evolução, e não foi feita tendo em vista que o ganho é pequeno em um núcleo de ordem única e sem execução especulativa, e a complexidade de verificação cresce de forma desproporcional ao escopo de uma monografia.

Por fim, a integração com o restante do núcleo condiciona o que o subsistema pode demonstrar. Enquanto não houver unidade de gerência de memória com tradução Sv32, interrupções de temporizador e entrada e saída mapeada em memória, o processador não executa um sistema operacional, e a avaliação fica restrita a programas sintéticos. São essas as etapas que separam o subsistema descrito aqui do objetivo coletivo que o motivou.
